\chapter{Health Checks}

\section*{Definition and Overview}
Health checks are routine medical assessments that look for early signs of disease, risk factors, and opportunities to improve health, in people who feel well or have no concerning symptoms. They include blood pressure and cholesterol measurement, blood glucose and kidney checks, cancer screening, and discussions about lifestyle, vaccination, and mental health. In Australia, the Medicare Benefits Schedule funds a comprehensive health assessment for people at certain ages and stages, including those aged 45--49 with risk factors and people aged 75 and over. Regular health checks detect problems early, when treatment is most effective, and provide an opportunity to prevent disease.

\section*{Epidemiology}
Many of the leading causes of illness and death in Australia are preventable or detectable early. Around half of adults have high blood pressure or elevated cholesterol, and many do not know it because these conditions cause no symptoms. Heart disease, stroke, type 2 diabetes, and common cancers account for a large share of deaths, and screening reduces their impact. Yet a significant proportion of Australians have not had a basic health check in the past two years. Regular checks are associated with earlier diagnosis, better risk factor control, and improved outcomes.

\section*{Aetiology and Pathophysiology}
Health checks work by identifying risk factors and disease at an early, treatable stage.

\begin{itemize}
    \item \textbf{Risk factor detection:} blood pressure, cholesterol, blood glucose, and weight identify people at risk of heart disease, stroke, and diabetes
    \item \textbf{Early disease detection:} screening finds cancers and other conditions before symptoms appear
    \item \textbf{Prevention:} vaccination and lifestyle advice prevent disease
    \item \textbf{Personalised assessment:} checks are tailored to age, sex, family history, and risk factors
    \item \textbf{Monitoring:} follow-up checks track risk factors and the effect of treatment
    \item \textbf{Population screening:} programs such as cervical, breast, and bowel screening operate nationally
    \item \textbf{Opportunistic checks:} blood pressure and other measurements taken during any consultation
\end{itemize}

\section*{Clinical Features}
Health checks address different things at different ages and for different risk groups.

\begin{itemize}
    \item Blood pressure, cholesterol, and blood glucose measurement
    \item Weight, height, and body mass index
    \item Smoking, alcohol, diet, and physical activity assessment
    \item Cancer screening: cervical, breast, bowel, and prostate discussions
    \item Heart disease risk assessment, including absolute cardiovascular risk
    \item Kidney function testing for people with diabetes or high blood pressure
    \item Bone health assessment for people at risk of osteoporosis
    \item Vaccination review and catch-up
    \item Mental health screening and discussion
    \item Vision, hearing, and dental checks as appropriate
\end{itemize}

\section*{Differential Diagnosis}
A health check is not a diagnosis in itself, but its findings prompt further assessment.

\begin{itemize}
    \item Elevated blood pressure requiring confirmation and review
    \item Elevated cholesterol or glucose requiring repeat testing and assessment
    \item Abnormal screening results requiring confirmatory tests
    \item Risk factors requiring lifestyle and, where appropriate, medicine treatment
    \item Mental health concerns requiring assessment
    \item Over-screening, which should be avoided when tests do more harm than good
    \item Under-screening in people who have not attended checks
\end{itemize}

\section*{When to Seek Medical Care}
Everyone should have regular health checks, with the frequency depending on age and risk. Medicare funds a health assessment for people aged 45--49 with risk factors, and annual assessments for people aged 75 and over. See a doctor sooner for any new or concerning symptoms.

\textbf{Call 000 (or local emergency services) for:}
\begin{itemize}
    \item Chest pain, breathlessness, or symptoms of stroke, which need emergency care rather than a routine check
    \item Sudden severe symptoms at any time
    \item Any medical emergency, regardless of recent health checks
\end{itemize}

\section*{Diagnosis and Investigations}
The components of a health check depend on the individual.

\begin{itemize}
    \item \textbf{History:} symptoms, lifestyle, family history, medicines, and vaccination status
    \item \textbf{Measurements:} blood pressure, weight, height, and BMI
    \item \textbf{Blood tests:} cholesterol, blood glucose, kidney function, and other tests based on risk
    \item \textbf{Heart disease risk assessment:} an absolute risk score from age, blood pressure, cholesterol, smoking, and diabetes
    \item \textbf{Cancer screening:} cervical screening, breast screening, and the bowel screening kit
    \item \textbf{Physical examination:} heart, lungs, abdomen, and other systems as indicated
    \item \textbf{Follow-up plan:} a clear plan for lifestyle changes, tests, and review
\end{itemize}

\section*{Management}
The value of a health check lies in acting on its findings.

\begin{itemize}
    \item \textbf{Lifestyle advice:} smoking cessation, healthy eating, physical activity, weight management, and alcohol reduction
    \item \textbf{Blood pressure management:} lifestyle measures and, if needed, medicines
    \item \textbf{Cholesterol and glucose management:} dietary change and medicines as indicated
    \item \textbf{Screening follow-up:} attending recommended cancer screening and acting on results
    \item \textbf{Vaccination:} completing recommended immunisations
    \item \textbf{Mental health:} advice and referral for any identified concerns
    \item \textbf{Medicines review:} ensuring medicines are appropriate and safe
    \item \textbf{Scheduled review:} a date for the next check or follow-up
\end{itemize}

\section*{Complications}
\begin{itemize}
    \item Missed diagnosis when checks are not attended
    \item Anxiety from abnormal results that are not properly explained
    \item Over-diagnosis and over-treatment from excessive testing
    \item False reassurance from normal results if risk factors change
    \item Unnecessary costs from tests not aligned with evidence
    \item Failure to act on findings, which negates the benefit
\end{itemize}

\section*{Prognosis and Outcomes}
Regular health checks improve outcomes by detecting risk factors and disease early. Blood pressure and cholesterol treatment reduces heart attacks and strokes, diabetes screening prevents complications, and cancer screening reduces deaths from cervical, breast, and bowel cancer. People who attend health checks are more likely to have their risk factors controlled and less likely to present with advanced disease. The greatest benefits come from attending checks consistently and acting on the results.

\section*{Prevention and Self-Care}
\begin{itemize}
    \item Book a health check with your GP, including the Medicare-funded assessment if you are eligible
    \item Know your blood pressure, cholesterol, and blood glucose numbers
    \item Complete the bowel screening kit when it arrives in the mail
    \item Attend cervical and breast screening according to the schedule
    \item Provide an accurate family history to your doctor
    \item Bring a list of your medicines to every check
    \item Follow up on any abnormal results and attend review appointments
    \item Maintain healthy habits between checks: diet, exercise, sleep, and no smoking
    \item See a doctor promptly for any new symptoms rather than waiting for the next check
\end{itemize}

\section*{Sources and Further Reading}
The following reputable sources provide further information for healthcare students and patients seeking reliable, current advice about health checks.

\begin{itemize}
    \item Healthdirect Australia. \textit{Health checks and assessments.} https://www.healthdirect.gov.au/
    \item Australian Government Department of Health and Aged Care. \textit{Medicare health assessments.} https://www.health.gov.au/
    \item Royal Australian College of General Practitioners. \textit{Preventive health and screening guidelines.} https://www.racgp.org.au/
    \item Australian Institute of Health and Welfare. \textit{Chronic disease risk factors.} https://www.aihw.gov.au/
    \item NHS. \textit{NHS Health Check.} https://www.nhs.uk/
    \item US Preventive Services Task Force. \textit{Screening recommendations.} https://www.uspreventiveservicestaskforce.org/
\end{itemize}
